% ============================================================
\setpartpreamble{\partintro
  {Betrieb in Produktion: Row Level Security, mandantensichere Beziehungen, Nachrüsten,
  Performance und der Weg vom Plugin in die Anwendung und in die CI.}
  {
    \item Kapitel~\ref{ch:row-level-security-in-production}: Row Level Security in
    Produktion
    \item Kapitel~\ref{ch:tenant-safe-relationships}: Mandantensichere Beziehungen
    \item Kapitel~\ref{ch:retrofitting-and-operating}: Mandantenfähigkeit nachrüsten und
    betreiben
    \item Kapitel~\ref{ch:plugin-performance}: Performance von Plugins
    \item Kapitel~\ref{ch:using-the-plugin-in-asp}: Einsatz in ASP.NET Core
    \item Kapitel~\ref{ch:ci-and-api-stability}: CI und API-Stabilität
  }}
\part{Betrieb und Integration}\label{part:operations}

% ------------------------------------------------------------
\chapter{Row Level Security in Produktion}\label{ch:row-level-security-in-production}
\versionrisk{stable}

RLS ist nur so stark wie seine Betriebsbedingungen. Diese Fallen machen die Policies
stillschweigend wirkungslos oder unbrauchbar:

\begin{datatable}{H{0.93}H{0.93}H{1.14}}
  \textbf{Falle} & \textbf{Folge} & \textbf{Abhilfe} \\
  \tablehead
  App läuft als Tabelleneigentümer & Policies werden ignoriert & \code{FORCE ROW LEVEL
  SECURITY} (Default im Plugin) und getrennte Rollen \\
  \tablerule
  Superuser oder \code{BYPASSRLS} & RLS wird immer umgangen, auch mit FORCE & App-Rolle mit
  \code{NOSUPERUSER NOBYPASSRLS} (eigene Operation,
  Kapitel~\ref{ch:custom-migration-operations}) \\
  \tablerule
  Setting nicht gesetzt & Ohne \code{nullif}: Fehler; mit: leere Ergebnisse & Interceptor
  setzt den Wert bei jedem Öffnen \\
  \tablerule
  Kein Index auf der Mandantenspalte & Jede Abfrage scannt die ganze Tabelle &
  Index-Convention (Kapitel~\ref{ch:conventions}) \\
  \tablerule
  Transaction-Pooling (PgBouncer) & Session-Settings landen bei anderen Clients oder gehen
  verloren & \code{Use\allowbreak{}Transaction\allowbreak{}Pooling(\allowbreak{})}: Setting
  transaktionslokal (siehe unten) \\
  \tablerule
  Verbindung außerhalb von EF geöffnet & Die Session behält einen alten Mandanten oder
  keinen & EF die Verbindungen öffnen lassen oder
  \code{Apply\allowbreak{}Tenant\allowbreak{}Session(\allowbreak{})}
  (Kapitel~\ref{ch:interceptors}) \\
  \tablerule
  „Bypass“-Schalter für die App-Rolle & Wer SQL als App-Rolle ausführen kann, kann ihn
  umlegen & Eigene Rollen für mandantenübergreifenden Zugriff (siehe unten) \\
\end{datatable}

\section{Das Bedrohungsmodell}\label{sec:threat-model}

RLS ist eine Schutzschicht mit klaren Grenzen. Wer sie einsetzt, sollte wissen, wovor sie
schützt, worauf sie sich verlässt und wo sie aufhört:

\begin{datatable}{H{0.45}H{1.55}}
  \textbf{Aspekt} & \textbf{Im Beispiel} \\
  \tablehead
  Geschützt & Die Zeilen mandantenbezogener Tabellen in den unterstützten Abbildungen
  (Seite~\pageref{sec:supported-table-mappings}). Nicht geschützt sind gemeinsame Stammdaten,
  das Schema und, ohne zusammengesetzte Schlüssel, die Existenz fremder Werte
  (Kapitel~\ref{ch:tenant-safe-relationships}). \\
  \tablerule
  Vertrauenswürdig & Der Datenbankserver und seine Rollen; die Migrationsrolle als
  Tabelleneigentümer; der Code von Anwendung und Plugin, also auch die Interceptors; der
  \code{ITenant\allowbreak{}Provider} und die Quelle, aus der er den Mandanten liest. \\
  \tablerule
  Annahmen & Die App-Rolle ist weder Superuser noch Eigentümer und hat kein
  \code{BYPASSRLS}. Der Mandant stammt aus einem authentifizierten Claim, nicht aus einem
  Header oder Parameter, den der Client frei wählt. Verbindungen werden von EF geöffnet, und es
  gibt kein Transaction-Pooling, außer mit
  \code{Use\allowbreak{}Transaction\allowbreak{}Pooling(\allowbreak{})}; dann kommen die
  Transaktionen von EF (Seite~\pageref{sec:transaction-pooling}). Das Modell nutzt nur
  Konfigurationen, die der Vertrag auf Seite~\pageref{sec:provider-comparison} als unterstützt
  ausweist. \\
  \tablerule
  Grenze & Die App-Rolle darf das Session-Setting selbst setzen; das muss sie, damit der
  Interceptor es setzen kann. Wer beliebiges SQL unter dieser Rolle ausführt, wechselt damit
  den Mandanten. \\
\end{datatable}

Die Grenze ist kein Fehler des Plugins, sondern eine Eigenschaft des Session-Modells.
\code{Threat\allowbreak{}Model\allowbreak{}Tests} zeigt sie gegen PostgreSQL: Ein
\code{SELECT set\_\allowbreak{}config(\allowbreak{}'app.\allowbreak{}tenant\_\allowbreak{}id',
\allowbreak{} .\allowbreak{}.\allowbreak{}.\allowbreak{})} unter der App-Rolle genügt, und die
Abfrage danach liefert die Zeilen des anderen Mandanten.
\code{Sql\allowbreak{}Server\allowbreak{}Threat\allowbreak{}Model\allowbreak{}Tests} zeigt
dasselbe mit \code{sp\_\allowbreak{}set\_\allowbreak{}session\_\allowbreak{}context}. Daraus
folgt, wovor RLS hier schützt und wovor nicht:

\begin{itemize}
  \item \textbf{Geschützt:} vergessene Filter, fehlerhafte LINQ-Abfragen, rohes SQL der
    Anwendung mit falschem \code{WHERE}, Massenoperationen, Berichte, die versehentlich über
    die App-Rolle laufen. In all diesen Fällen setzt das Plugin den Mandanten, und der Code
    ändert ihn nicht.
  \item \textbf{Nicht geschützt:} SQL-Injection und kompromittierter Anwendungscode. Beides
    kann das Setting ändern, auch innerhalb eines einzigen Statements, weil
    \code{set\_\allowbreak{}config} eine Funktion ist. Dagegen helfen Parameter statt
    zusammengesetzter Strings (EF parametrisiert LINQ und
    \code{From\allowbreak{}Sql}/\code{Sql\allowbreak{}Query} mit interpolierten Strings;
    \code{From\allowbreak{}Sql\allowbreak{}Raw} mit verketteten Werten nicht) und Rechte, die
    der App-Rolle nur die Tabellen geben, die sie braucht.
  \item \textbf{Außerhalb der Datenbank:} Ein gefälschter oder falsch zugeordneter Claim gibt
    Zugriff auf genau diesen Mandanten. Das ist Sache der Authentifizierung; das Plugin
    übernimmt den Mandanten aus dem \code{ITenant\allowbreak{}Provider}.
\end{itemize}

Mögliche Härtungen setzt das Beispiel bewusst nicht um. SQL Server kann den Schlüssel mit
\code{sp\_\allowbreak{}set\_\allowbreak{}session\_\allowbreak{}context
.\allowbreak{}.\allowbreak{}.\allowbreak{}, @read\_only = 1} bis zum Ende der Verbindung
einfrieren
(\href{https://learn.microsoft.com/en-us/sql/relational-databases/security/row-level-security}{SQL Server: Row-Level Security});
der Test oben belegt, dass ein späterer Wechsel dann scheitert. Der Preis: Auch
\code{Apply\allowbreak{}Tenant\allowbreak{}Session} kann den Mandanten einer offenen Verbindung
nicht mehr ändern. PostgreSQL kennt kein schreibgeschütztes Setting; dort bleiben eine Rolle pro
Mandant mit \code{SET ROLE} oder ein signiertes Mandanten-Token, das eine
\code{SECURITY DEFINER}-Funktion prüft. Beides verlagert Komplexität in die Datenbank und ist
hier nur als Richtung genannt. Mandantenübergreifender Zugriff gehört weiterhin in eigene
Rollen (Abschnitt unten).

\section{Transaction-Pooling}\label{sec:transaction-pooling}

PgBouncer im Modus „transaction“ teilt eine Serververbindung zwischen Clients pro
Transaktion. Ein mit
\code{set\_\allowbreak{}config(\allowbreak{}.\allowbreak{}.\allowbreak{}.\allowbreak{},\allowbreak{}
false)} gesetzter Session-Wert kann dann bei einem anderen Client ankommen, eine
Sicherheitslücke. Die sichere Variante setzt den Wert mit
\code{set\_\allowbreak{}config(\allowbreak{}.\allowbreak{}.\allowbreak{}.\allowbreak{},\allowbreak{}
true)} \emph{transaktionslokal}: Er endet mit Commit oder Rollback. Dafür wechselt der
Interceptor von \code{IDb\allowbreak{}Connection\allowbreak{}Interceptor} zu
\code{IDb\allowbreak{}Transaction\allowbreak{}Interceptor} und setzt den Mandanten in
\code{Transaction\allowbreak{}Started}, im Befehl der neuen Transaktion. Das Beispiel bietet das
als Option an:

\begin{sourcecode}[Anwendung]{csharp}
services.AddDbContext<ShopContext>((sp, options) => options
    .UseNpgsql(connectionString)
    .UseAcmeTenancy(sp, t => t.UseTransactionPooling()));
\end{sourcecode}

% source: src/Acme.EntityFrameworkCore.Tenancy.PostgreSQL/Runtime/PostgreSqlTransactionSessionInterceptor.cs
\begin{sourcecode}[Runtime/\allowbreak{}Postgre\allowbreak{}Sql\allowbreak{}Transaction\allowbreak{}Session\allowbreak{}Interceptor.\allowbreak{}cs (\allowbreak{}gekürzt)]{csharp}
public override async ValueTask<DbTransaction> TransactionStartedAsync(
    DbConnection connection,
    TransactionEndEventData eventData,
    DbTransaction result,
    CancellationToken cancellationToken = default)
{
    await using var command = CreateSessionCommand(connection, result, eventData.Context);
    await command.ExecuteNonQueryAsync(cancellationToken);
    return result;
}

// SaveChanges immer in einer Transaktion, auch bei einem einzelnen Befehl
private static void RequireTransaction(DbContextEventData eventData)
{
    // ...
        context.Database.AutoTransactionBehavior = AutoTransactionBehavior.Always;
}

protected virtual DbCommand CreateSessionCommand(
    DbConnection connection, DbTransaction transaction, DbContext? context)
{
    // ...
    var command = connection.CreateCommand();
    command.Transaction = transaction;
    command.CommandText = "SELECT set_config(@name, @value, true)";   // is_local = true
    // ...
}
\end{sourcecode}

\code{Use\allowbreak{}Acme\allowbreak{}Tenancy(\allowbreak{}sp)} registriert dann diesen
Interceptor statt des Connection-Interceptors. Drei Regeln gehören dazu:

\begin{itemize}
  \item \textbf{Lesen in Transaktionen.} Abfragen außerhalb einer Transaktion haben keinen
    Mandanten und sehen keine Zeilen. Das ist sicher, aber die Anwendung muss lesende Arbeit,
    die Daten braucht, in eine Transaktion legen.
  \item \textbf{Speichern immer in einer Transaktion.} EF lässt die automatische Transaktion
    weg, wenn \code{Save\allowbreak{}Changes} nur einen Befehl ausführt
    (\code{Auto\allowbreak{}Transaction\allowbreak{}Behavior.\allowbreak{}When\allowbreak{}Needed}).
    Ein einzelner Insert liefe dann ohne Mandanten und scheiterte an der Policy. Der Interceptor
    ist deshalb zugleich ein \code{ISave\allowbreak{}Changes\allowbreak{}Interceptor} und stellt
    \code{Always} ein.
  \item \textbf{Kein Session-Wert.} Ist die Option gesetzt und trotzdem ein
    \code{Tenant\allowbreak{}Connection\allowbreak{}Interceptor} registriert, verweigert dieser
    das Setzen mit einer \code{Tenancy\allowbreak{}Exception}, statt einen Wert zu
    hinterlassen, der die Transaktion überdauert. Damit entfällt auch
    \code{Apply\allowbreak{}Tenant\allowbreak{}Session}: Ein Aufruf scheitert mit dem Hinweis,
    dass der Mandant pro Transaktion gesetzt wird und die Befehle in eine gehören.
\end{itemize}

Der Interceptor wirkt nur auf Transaktionen, die EF selbst beginnt. Damit gehört zum Vertrag,
auf welchem Weg eine Transaktion entsteht. Die mittlere Spalte gilt auch für SQL Server, das
keinen Transaction-Pooling-Modus hat:

\begin{datatable}{H{0.80}H{1.10}H{1.10}}
  \textbf{Transaktionsweg} & \textbf{Session-Modus (Standard)} &
  \textbf{\code{Use\allowbreak{}Transaction\allowbreak{}Pooling(\allowbreak{})}} \\
  \tablehead
  \code{Begin\allowbreak{}Transaction(\allowbreak{})} über EF & Unterstützt und getestet: EF
  öffnet die Verbindung, der Interceptor setzt die Session & Unterstützt und getestet: Mandant
  in \code{Transaction\allowbreak{}Started}, nach Commit und Rollback weg \\
  \tablerule
  \code{Save\allowbreak{}Changes} ohne eigene Transaktion & Unterstützt und getestet &
  Unterstützt und getestet: Der Interceptor erzwingt eine Transaktion, auch für einen einzelnen
  Befehl \\
  \tablerule
  Abfrage außerhalb einer Transaktion & Unterstützt und getestet & Kein Mandant, keine Zeilen
  (fail closed, getestet) \\
  \tablerule
  \code{Use\allowbreak{}Transaction(\allowbreak{})} mit eigener Transaktion & Nur unter
  Voraussetzungen: EF hat die Verbindung geöffnet, oder
  \code{Apply\allowbreak{}Tenant\allowbreak{}Session(\allowbreak{})} nach dem Öffnen außerhalb
  von EF (\code{Connection\allowbreak{}Lifetime\allowbreak{}Tests}); der Aufruf selbst ist nicht
  verifiziert & Nicht unterstützt: EF meldet \code{Transaction\allowbreak{}Used}, nicht
  \code{Transaction\allowbreak{}Started}; kein Mandant, keine Zeilen (fail closed, getestet) \\
  \tablerule
  Umgebende \code{Transaction\allowbreak{}Scope} & Unterstützt und getestet: Die Verbindung
  tritt beim Öffnen bei & Nicht unterstützt: EF beginnt keine eigene Transaktion; Lesen sieht
  keine Zeilen, Schreiben scheitert an der Policy (fail closed, getestet) \\
\end{datatable}

\code{Transaction\allowbreak{}Pooling\allowbreak{}Tests} prüft gegen PostgreSQL: Innerhalb
der Transaktion ist der Wert gesetzt und RLS greift, nach Commit und nach Rollback ist er auf
derselben Verbindung leer, ein einzelner Insert bekommt seine Transaktion, der
Connection-Interceptor lehnt den Session-Wert ab,
\code{Apply\allowbreak{}Tenant\allowbreak{}Session} erklärt, warum es nicht verfügbar ist, und
\code{Use\allowbreak{}Transaction} sowie \code{Transaction\allowbreak{}Scope} scheitern
geschlossen. PgBouncer selbst ist nicht Teil des Tests;
geprüft wird die Eigenschaft, auf die es ankommt: dass nach der Transaktion kein Wert bleibt.

\section{Performance von Policies}

PostgreSQL fügt Policies wie zusätzliche \code{WHERE}-Bedingungen in jede Abfrage ein.
\code{current\_\allowbreak{}setting} ist \code{STABLE}. Das garantiert innerhalb eines
Statements dasselbe Ergebnis für alle Zeilen und \emph{erlaubt} dem Optimierer, mehrere
Aufrufe zu einem zusammenzufassen
(\href{https://www.postgresql.org/docs/17/xfunc-volatility.html}{PostgreSQL: Function Volatility Categories}); eine
feste Zahl von Aufrufen garantiert es nicht. Wie oft die Funktion tatsächlich läuft, hängt
vom Plan ab. Steht der Vergleich in einer \code{Index Cond}, wertet PostgreSQL den
Vergleichswert einmal aus und liest nur die Zeilen des Mandanten. Steht er als
\code{Filter} über einem Seq Scan, kann er pro Zeile ausgewertet werden, und dann wiegt
der Scan schwerer als die Funktion. Teure Policies entstehen durch Unterabfragen (etwa
Rechte aus einer anderen Tabelle). Zusätzlich hält PostgreSQL Bedingungen aus der Abfrage,
deren Funktionen nicht \code{LEAKPROOF} sind, hinter den Policy-Bedingungen zurück, was
Optimierungen verhindern kann (\href{https://www.postgresql.org/docs/17/ddl-rowsecurity.html}{PostgreSQL: Row Security Policies}).

Aussagekräftig ist deshalb nur der Plan unter realen Bedingungen: \code{EXPLAIN
(\allowbreak{}ANALYZE,\allowbreak{} BUFFERS)} als App-Rolle mit gesetztem Mandanten, mit
realistischer Datenmenge und aktuellen Statistiken. Als Eigentümer oder Superuser gemessen fehlen
die Policies im Plan, weil Superuser und Rollen mit \code{BYPASSRLS} sie umgehen und der
Eigentümer ohne \code{FORCE} ebenso. Den Overhead zeigt der Vergleich derselben Abfrage mit
und ohne Policy, nicht die Zahl der Funktionsaufrufe.

\section{Legitimer mandantenübergreifender Zugriff}

Früher oder später braucht jemand Daten über Mandanten hinweg: Berichte für den Betreiber,
Wartungsjobs, Support. Die verlockende Abkürzung ist ein Schalter in der Anwendung, etwa ein
zusätzliches Session-Setting \code{app.\allowbreak{}bypass}, das die Policy berücksichtigt.
Das ist keine Sicherheitsgrenze: Was die Anwendungsrolle setzen kann, kann auch ein
eingeschleustes SQL-Statement setzen. Die Trennung muss stattdessen über Datenbankrollen
entstehen:

\begin{datatable}{H{0.75}H{1.25}H{1.00}}
  \textbf{Bedarf} & \textbf{Lösung} & \textbf{Warum es sicher ist} \\
  \tablehead
  Berichte, Analysen & Eigene Nur-Lese-Rolle mit \code{BYPASSRLS}, nur \code{SELECT},
  eigenem Connection-String und einer Context-Konfiguration ohne Mandant & Die App-Rolle
  kann nicht zu dieser Rolle werden; ihre Zugangsdaten liegen woanders \\
  \tablerule
  Support liest ausgewählte Daten & Zusätzliche permissive Policy für eine Support-Rolle:
  \code{Has\allowbreak{}Policy(\allowbreak{}.\allowbreak{}.\allowbreak{}.\allowbreak{})} mit
  \code{For(\allowbreak{}Select)} und
  \code{To(\allowbreak{}\textquotedbl{}shop\_\allowbreak{}support\textquotedbl{})} &
  Beschränkt auf SELECT und auf Mitglieder der Rolle \\
  \tablerule
  Wartung pro Mandant & Hintergrundjob, der die Mandanten durchläuft und jeden einzeln setzt
  (Kapitel~\ref{ch:using-the-plugin-in-asp}) & Läuft durch die normalen Policies \\
\end{datatable}

\begin{sourcecode}[Reporting-Rolle (\allowbreak{}tests/\allowbreak{}.\allowbreak{}.\allowbreak{}.\allowbreak{}Integration\allowbreak{}Tests/\allowbreak{}Cross\allowbreak{}Tenant\allowbreak{}Access\allowbreak{}Tests.\allowbreak{}cs,\allowbreak{} vereinfacht)]{sql}
-- Angelegt von einem Administrator: BYPASSRLS erfordert einen Superuser
CREATE ROLE shop_reporting LOGIN PASSWORD '...' NOSUPERUSER BYPASSRLS;
GRANT USAGE ON SCHEMA public TO shop_reporting;
GRANT SELECT ON ALL TABLES IN SCHEMA public TO shop_reporting;
\end{sourcecode}

Die Anwendungsrolle dagegen kann die Policies nicht abschalten. \code{SET
row\_\allowbreak{}security = off} deaktiviert RLS für eine Rolle ohne \code{BYPASSRLS}
nicht; PostgreSQL lehnt stattdessen jede Abfrage ab, die die Policies betreffen würden.
\code{Cross\allowbreak{}Tenant\allowbreak{}Access\allowbreak{}Tests} prüft alle drei Punkte:
Die Reporting-Rolle sieht alle Mandanten, sie kann nicht schreiben, und die App-Rolle kann
sich nicht selbst helfen.

\begin{important}[Mandantenübergreifenden Zugriff protokollieren]
  Jede Verbindung der Reporting-Rolle ist ein Zugriff auf alle Mandanten. Sie gehört
  protokolliert (\code{log\_\allowbreak{}connections}, pgaudit), und die Rolle sollte so
  wenigen Personen und Diensten wie möglich zugewiesen werden.
  \code{Create\allowbreak{}Application\allowbreak{}Role} des Plugins legt Rollen bewusst
  immer mit \code{NOBYPASSRLS} an: Bypass-Rollen gehören nicht in die Migrationen der
  Anwendung.
\end{important}

% ------------------------------------------------------------
\chapter{Mandantensichere Beziehungen}\label{ch:tenant-safe-relationships}
\versionrisk{stable}

Welche Tabellenabbildungen das Plugin schützt und welche es beim Modellbau ablehnt, legt die
Matrix auf Seite~\pageref{sec:supported-table-mappings} fest; dieses Kapitel wiederholt sie
nicht. Es geht um die zweite Achse: Beziehungen zwischen geschützten Tabellen. RLS regelt, welche
\emph{Zeilen} eine Sitzung sieht. Ob eine Zeile nur auf Zeilen desselben Mandanten verweist,
regelt RLS nicht. Das ist \emph{referenzielle Mandantenkonsistenz}, und sie entsteht im Modell.

Das Beispiel verwendet Kunde, Bestellung und Position: \code{Customer} → \code{Order} →
\code{Order\allowbreak{}Item}, alle drei \code{ITenant\allowbreak{}Owned}, dazu ein gemeinsames
\code{Product} ohne Mandant.

\section{Sichtbarkeit ist nicht Konsistenz}

PostgreSQL dokumentiert ausdrücklich, dass Prüfungen von Fremdschlüsseln, eindeutigen und
Primärschlüsseln RLS umgehen, damit die Datenintegrität erhalten bleibt, und warnt vor
„covert channels“ durch diese Prüfungen
(\href{https://www.postgresql.org/docs/17/ddl-rowsecurity.html}{PostgreSQL: Row Security Policies}).
Die SQL-Server-Dokumentation sagt dazu nichts ausdrücklich;
\code{Sql\allowbreak{}Server\allowbreak{}Threat\allowbreak{}Model\allowbreak{}Tests} zeigt,
dass die Fremdschlüsselprüfung dort ebenfalls Zeilen sieht, die das Filter-Prädikat verbirgt.
Mit Fremdschlüsseln nur auf der ID, wie EF sie standardmäßig anlegt, folgt daraus:

\begin{datatable}{H{0.70}H{1.30}}
  \textbf{Beobachtung} & \textbf{Ursache} \\
  \tablehead
  Mandant A legt eine Position zu einer Bestellung von B an & Die FK-Prüfung findet die
  Bestellung, obwohl A sie nicht sieht. \\
  \tablerule
  A erfährt, ob eine fremde ID existiert & Unbekannte IDs scheitern am Fremdschlüssel,
  fremde nicht. \\
  \tablerule
  Die eigene Position verschwindet aus \code{Include(\allowbreak{}i => i.\allowbreak{}Order)} &
  Eine Pflichtnavigation wird ein \code{INNER JOIN}, und die Bestellung ist unsichtbar. \\
  \tablerule
  A löscht eine eigene Bestellung, eine Position von B verschwindet & Die kaskadierende
  Löschung ist eine referenzielle Aktion und umgeht RLS ebenfalls. \\
  \tablerule
  A erfährt, dass eine E-Mail-Adresse bei B existiert & Ein globaler Unique-Index meldet die
  Kollision. \\
\end{datatable}

\code{Relationship\allowbreak{}Tests} weist jede Zeile gegen PostgreSQL mit eingeschränkter
Rolle nach. Keiner dieser Fälle ist ein Leseleck im engeren Sinn; A sieht keine fremde Zeile.
Aber die Daten verlieren ihre Mandantengrenze, und die Fehlermeldungen beantworten Fragen, die
A nicht stellen dürfte.

\section{Zusammengesetzte Fremdschlüssel}

Die Lösung ist, den Mandanten zum Teil jeder Referenz zu machen. Der Prinzipal bekommt einen
alternativen Schlüssel aus Mandant und ID, die abhängige Seite einen Fremdschlüssel aus ihrer
Mandantenspalte und der ID:

% source: tests/Acme.EntityFrameworkCore.Tenancy.IntegrationTests/RelationshipTests.cs
\begin{sourcecode}[tests/\allowbreak{}.\allowbreak{}.\allowbreak{}.\allowbreak{}Integration\allowbreak{}Tests/\allowbreak{}Relationship\allowbreak{}Tests.\allowbreak{}cs › Tenant\allowbreak{}Safe\allowbreak{}Context]{csharp}
modelBuilder.Entity<Customer>(b =>
{
    b.HasAlternateKey(c => new { c.TenantId, c.Id });
    b.HasIndex(c => new { c.TenantId, c.Email }).IsUnique();   // eindeutig pro Mandant
});

modelBuilder.Entity<Order>(b =>
{
    b.HasAlternateKey(o => new { o.TenantId, o.Id });
    b.HasOne(o => o.Customer).WithMany(c => c.Orders)
        .HasForeignKey(o => new { o.TenantId, o.CustomerId })
        .HasPrincipalKey(c => new { c.TenantId, c.Id })
        .OnDelete(DeleteBehavior.Restrict);                   // Kunden nicht kaskadieren
});

modelBuilder.Entity<OrderItem>()
    .HasOne(i => i.Order).WithMany(o => o.Items)
    .HasForeignKey(i => new { i.TenantId, i.OrderId })
    .HasPrincipalKey(o => new { o.TenantId, o.Id })
    .OnDelete(DeleteBehavior.Cascade);                         // Positionen gehören dazu
\end{sourcecode}

Jetzt verweist eine Position von A auf das Paar (A, ID). Eine Bestellung von B erfüllt das nie,
und der Fehler ist derselbe wie für eine unbekannte ID: Es gibt nichts mehr zu erfahren. Die
Mandantenspalte ist zugleich Teil des Fremdschlüssels und Mandanten-Property. Der
SaveChanges-Interceptor setzt sie bei neuen Positionen, und das Fixup über die Navigation setzt
denselben Wert. Ein Index für die Policy ist nicht nötig, weil Schlüssel und Fremdschlüssel mit
der Mandantenspalte beginnen; die Index-Convention erkennt das.

Der Preis ist ein zusätzlicher eindeutiger Index pro Prinzipal und eine breitere
Fremdschlüsselspalte. Das Plugin legt diese Schlüssel nicht automatisch an. Eine Convention
könnte es, müsste aber jede Beziehung zwischen geschützten Typen umschreiben und mit expliziter
Konfiguration abgleichen. Das Beispiel überlässt die Entscheidung dem Modell der Anwendung.

\section{Eindeutigkeit pro Mandant}

Ein Unique-Index über eine fachliche Spalte gilt für die ganze Tabelle, also über alle
Mandanten. Fast immer ist Eindeutigkeit pro Mandant gemeint: \code{(TenantId, Email)} statt
\code{Email}. Wo Eindeutigkeit wirklich global sein muss, etwa bei Subdomains der Mandanten,
gehört die Spalte in das Mandantenverzeichnis, das bewusst keine Mandantendaten enthält.

\section{Gemeinsame Stammdaten}

\code{Product} hat keinen Mandanten, keine Policy und ist für alle lesbar. Fremdschlüssel von
Mandantendaten auf Stammdaten sind unproblematisch; der Test liest das Produkt über die Position.
Die Gegenrichtung, Stammdaten mit Verweisen auf Mandantendaten, darf es nicht geben. Schreiben
in Stammdaten sollte nicht die App-Rolle dürfen; \code{Grant\allowbreak{}Table\allowbreak{}Privileges}
gewährt Rechte auf alle Tabellen eines Schemas, eine feinere Vergabe ist Sache der Migrationen
der Anwendung.

\section{Join-Tabellen}

Implizite Many-to-many-Tabellen lehnt das Plugin ab, wenn eine Seite mandantenbezogen ist. Eine
Join-Entität mit eigener Mandanten-Property ist geschützt
(Seite~\pageref{sec:supported-table-mappings}). Für Konsistenz bekommt sie dieselben
zusammengesetzten Fremdschlüssel wie jede abhängige Tabelle: \code{(TenantId, LeftId)} und
\code{(TenantId, RightId)} auf die alternativen Schlüssel beider Seiten. Dann kann eine
Verknüpfung nur zwei Zeilen desselben Mandanten verbinden.

\section{Löschregeln}

Kaskadierende Löschungen laufen in der Datenbank als referenzielle Aktionen und umgehen RLS.
Mit zusammengesetzten Fremdschlüsseln bleiben sie trotzdem im Mandanten, weil jede abhängige
Zeile denselben Mandanten trägt. Ohne sie löscht eine Kaskade fremde Zeilen; der Test
\code{Naive\allowbreak{}\_cascade\allowbreak{}\_deletes\allowbreak{}\_rows\allowbreak{}\_of\allowbreak{}\_another\allowbreak{}\_tenant}
zeigt genau das. Für Wurzeln wie den Kunden ist \code{Restrict} meist richtig, für Teile eines
Aggregats wie Positionen \code{Cascade}. EF kaskadiert zusätzlich im Change Tracker, aber nur
für geladene Entitäten; maßgeblich ist die Regel in der Datenbank.


% ------------------------------------------------------------
\chapter{Mandantenfähigkeit nachrüsten und betreiben}\label{ch:retrofitting-and-operating}
\versionrisk{stable}

Die bisherigen Kapitel gehen von einer neuen Datenbank aus. Meist gibt es aber schon Tabellen
mit Daten, eine laufende Anwendung und Nutzer, die während der Umstellung weiterarbeiten. Dieses
Kapitel zeigt den Weg von einer Tabelle ohne Mandantenspalte zu RLS und danach den Betrieb unter
Retries und Verbindungsabbrüchen. Alle Schritte stehen ausführbar in
\code{Operations\allowbreak{}Tests} gegen PostgreSQL. Den Abschluss bildet das Deployment:
Migrationen in der Pipeline mit getrennten Rollen und Änderungen ohne Ausfallzeit.

\section{Von der Tabelle ohne Mandant zu RLS}

Wer eine bestehende Entität einfach \code{ITenant\allowbreak{}Owned} implementieren lässt,
bekommt eine einzige Migration: Spalte mit dem leeren Mandanten als Default, RLS, Policy, Index.
Danach sieht niemand mehr die alten Zeilen. Auch ein Backfill in derselben Migration hilft nicht
sicher: Mit \code{FORCE ROW LEVEL SECURITY} gilt die Policy auch für die Migrationsrolle, und
ohne gesetzten Mandanten findet ihr \code{UPDATE} keine Zeile. Der Weg führt deshalb über zwei
Migrationen, nach dem Muster Erweitern, Befüllen, Absichern:

\begin{datatable}{H{0.22}H{0.90}H{1.88}}
  \textbf{Nr.} & \textbf{Schritt} & \textbf{Vorgehen} \\
  \tablehead
  1 & Erweitern & Die Entität implementiert \code{ITenant\allowbreak{}Owned}, bleibt aber mit
  \code{Has\allowbreak{}No\allowbreak{}Row\allowbreak{}Level\allowbreak{}Security(\allowbreak{})}
  ausgenommen. Die Migration legt nur die Spalte an. \\
  \tablerule
  2 & Befüllen & Als Eigentümer, in derselben oder einer eigenen Migration: Jede Zeile bekommt
  ihren Mandanten aus dem, was die Zuordnung kennt, etwa einer Zuordnungstabelle oder dem
  Mandanten des Eltern-Datensatzes. \\
  \tablerule
  3 & Anwendung umstellen & Eine App-Version, deren Modell die Ausnahme nicht mehr hat:
  Der Interceptor setzt den Mandanten bei neuen Zeilen. Die Datenbank schützt noch nicht, das
  Modell prüft das nicht. \\
  \tablerule
  4 & Lücke schließen & Zeilen, die zwischen Schritt 1 und 3 ohne Mandant entstanden sind,
  nachträglich zuordnen. \\
  \tablerule
  5 & Absichern & Zweite Migration: Ausnahme entfernt, also \code{ENABLE}/\code{FORCE},
  Policy und Index. Als Erstes ein Wächter, der abbricht, solange eine Zeile keinen Mandanten
  hat. \\
  \tablerule
  6 & Verschärfen & Optional zusammengesetzte Schlüssel und Unique-Indizes pro Mandant
  (Kapitel~\ref{ch:tenant-safe-relationships}). \\
\end{datatable}

Der Wächter ist ein eigener Schritt am Anfang von \code{Up}, vor den Operationen, die EF erzeugt:

\begin{sourcecode}[Wächter (\allowbreak{}SQL)]{postgresql}
DO $$ BEGIN
  IF EXISTS (SELECT 1 FROM "Tickets"
             WHERE "TenantId" = '00000000-0000-0000-0000-000000000000') THEN
    RAISE EXCEPTION 'Tickets without a tenant: backfill first';
  END IF;
END $$;
\end{sourcecode}

\begin{sourcecode}[Migrations/\allowbreak{}…\allowbreak{}\_Enable\allowbreak{}Ticket\allowbreak{}Isolation.\allowbreak{}cs (Beispiel)]{csharp}
protected override void Up(MigrationBuilder migrationBuilder)
{
    // Vor allem anderen: ohne vollständige Zuordnung nicht absichern
    migrationBuilder.Sql(TicketGuardSql);      // der Wächter oben, als Konstante

    // ... von EF erzeugt: ENABLE/FORCE ROW LEVEL SECURITY, CREATE POLICY, CREATE INDEX
}
\end{sourcecode}

Die Migration läuft in einer Transaktion. Schlägt der Wächter an, bleibt die Tabelle ohne
RLS, und kein halber Zustand entsteht;
\code{Guard\allowbreak{}\_stops\allowbreak{}\_enabling\allowbreak{}\_rls\allowbreak{}\_while\allowbreak{}\_rows\allowbreak{}\_lack\allowbreak{}\_a\allowbreak{}\_tenant}
prüft das. Nach vollständiger Zuordnung sieht jeder Mandant genau seine Zeilen.

\subsection*{App-Versionen, Down und Rollback}

Während der Umstellung laufen oft zwei App-Versionen gleichzeitig, und jede trifft auf einen der
beiden Datenbankzustände:

\begin{datatable}{H{0.95}H{1.00}H{1.05}}
  \textbf{Datenbank} & \textbf{App ohne Mandant} & \textbf{App mit Mandant} \\
  \tablehead
  Nach Migration 1 (Spalte, kein RLS) & Läuft; neue Zeilen ohne Mandant (Schritt 4) & Läuft,
  setzt den Mandanten; noch ohne Schutz \\
  \tablerule
  Nach Migration 2 (RLS) & Fail closed: liest nichts, Schreiben scheitert & Geschützt \\
\end{datatable}

Die alte App-Version darf also nicht mehr laufen, wenn Migration 2 eingespielt ist. Ein Rollback
der Anwendung auf eine Version ohne Mandant setzt deshalb erst das Down von Migration 2 voraus.
Das Down entfernt Policy und RLS; Spalte und Werte bleiben, die Tabelle ist aber wieder
ungeschützt. Der Test prüft beides. Ein Down von Migration 1 löscht die Spalte und damit die
Zuordnung; es gehört nur in die Entwicklung.

Policy-Änderungen folgen demselben Muster. Ändert sich der Setting-Name
(\code{Use\allowbreak{}Setting\allowbreak{}Name}), ersetzt die Migration die Policy in ihrer
Transaktion. Zwischen Migration und neuer App-Version liest die Policy ein Setting, das die
laufende App nicht setzt, und die Datenbank verweigert alles, bis beide zusammenpassen. Ein
Übergang, bei dem der Interceptor beide Settings setzt, ist denkbar, aber nicht implementiert.
Für SQL Server beschreibt Seite~\pageref{sec:sqlserver-session-settings} denselben Fall.

\section{Retries, Transaktionen und Verbindungsabbrüche}

Das Plugin setzt den Mandanten, wenn EF eine Verbindung öffnet. Daraus ergibt sich das Verhalten
unter den Resilienzmechanismen von EF:

\begin{itemize}
  \item \textbf{Retries.} Eine Execution Strategy wie \code{Enable\allowbreak{}Retry\allowbreak{}On\allowbreak{}Failure}
    wiederholt eine Operation nach einem transienten Fehler auf einer neu geöffneten Verbindung.
    Der Interceptor läuft dabei wieder, der Mandant ist gesetzt. Der Test beendet die Sitzung
    serverseitig mit \code{pg\_\allowbreak{}terminate\_\allowbreak{}backend}; die nächste Abfrage
    scheitert, wird wiederholt und liefert nur Zeilen des Mandanten.
  \item \textbf{Explizite Transaktionen.} Mit einer wiederholenden Strategie lehnt EF Befehle in
    einer selbst begonnenen Transaktion ab, weil ein Retry die ganze Transaktion wiederholen
    müsste. Die Transaktion gehört deshalb in
    \code{Create\allowbreak{}Execution\allowbreak{}Strategy(\allowbreak{}).\allowbreak{}Execute\allowbreak{}Async(\allowbreak{}…)}.
    Innerhalb öffnet EF die Verbindung einmal, setzt den Mandanten und hält sie bis zum Commit.
    Bricht die Verbindung mitten in der Transaktion ab, wiederholt die Strategie den ganzen
    Delegaten auf einer neuen Verbindung; Seiteneffekte außerhalb der Datenbank müssen dann
    idempotent sein.
  \item \textbf{Neuinitialisierung der Session.} Npgsql setzt den Zustand einer Verbindung
    zurück, wenn sie in den Pool zurückgeht, solange \code{No Reset On Close} nicht aktiviert
    ist (\href{https://www.npgsql.org/doc/connection-string-parameters.html}{Npgsql: Connection String Parameters}).
    SqlClient setzt den Session-Kontext ebenfalls zurück, wenn es eine Verbindung aus dem Pool
    wiederverwendet (\code{sp\_\allowbreak{}reset\_\allowbreak{}connection}). Das Plugin verlässt sich
    auf keines von beidem, sondern setzt den Mandanten bei jedem Öffnen; nur Verbindungen, die
    nicht EF öffnet, brauchen
    \code{Apply\allowbreak{}Tenant\allowbreak{}Session} (Kapitel~\ref{ch:interceptors}).
\end{itemize}

Für Transaction-Pooling setzt \code{Use\allowbreak{}Transaction\allowbreak{}Pooling(\allowbreak{})}
den Mandanten pro Transaktion (Seite~\pageref{sec:transaction-pooling}). Nicht implementiert ist
ein schreibgeschützter Session-Schlüssel (Seite~\pageref{sec:threat-model}).
\code{Operations\allowbreak{}Tests} enthält den Verbindungsabbruch und die Transaktion in der
Strategie; die Fälle offener Verbindungen stehen in
\code{Connection\allowbreak{}Lifetime\allowbreak{}Tests}.


\section{Deployment: Pipeline, Rollen und Zero-Downtime}\label{sec:deployment}

Die Beispiele legen die Datenbank in der Entwicklung mit
\code{Development\allowbreak{}Database} an. In Produktion gehört jede Schemaänderung in eine
Pipeline, die zwei Rollen strikt trennt:

\begin{itemize}
  \item \textbf{\code{shop\_\allowbreak{}owner}} besitzt Schema, Tabellen und Policies und führt
    die Migrationen aus.
  \item \textbf{\code{shop\_\allowbreak{}app}} ist die Rolle der Anwendung: nur
    \code{SELECT}, \code{INSERT}, \code{UPDATE}, \code{DELETE}, immer unter den Policies.
\end{itemize}

Die Anwendung ruft deshalb niemals
\code{db.\allowbreak{}Database.\allowbreak{}Migrate(\allowbreak{})} beim Start auf. Dafür
bräuchte ihre Rolle DDL-Rechte, und wer eine Tabelle ändern darf, darf auch
\code{DISABLE ROW LEVEL SECURITY} ausführen oder die Policy löschen. \code{FORCE} bindet den
Eigentümer an die Policy, schützt sie aber nicht vor ihm. Eine Lücke in der Anwendung, etwa
eine SQL-Injection, wäre dann nicht mehr durch die Datenbank begrenzt. Hinzu kommt der Betrieb:
Bei einem Rolling Deployment starten mehrere Instanzen gleichzeitig, und eine fehlgeschlagene
Migration soll die Auslieferung stoppen, nicht die Instanzen in einer Neustartschleife halten.

\subsection*{Einmalig: die Rollen}

Bevor die erste Pipeline läuft, legt die Datenbankadministration die Rollen an. Das ist keine
Migration, denn Migrationen laufen bereits als Eigentümer. Das Skript liegt im Repository:

% source: docs/deployment/setup-roles.sql
\begin{sourcecode}[docs/\allowbreak{}deployment/\allowbreak{}setup-\allowbreak{}roles.\allowbreak{}sql]{postgresql}
CREATE ROLE shop_owner LOGIN NOSUPERUSER NOBYPASSRLS PASSWORD :'owner_password';
CREATE ROLE shop_app LOGIN NOSUPERUSER NOBYPASSRLS PASSWORD :'app_password';
CREATE DATABASE shop OWNER shop_owner;

...
GRANT USAGE ON SCHEMA public TO shop_app;

ALTER DEFAULT PRIVILEGES FOR ROLE shop_owner IN SCHEMA public
    GRANT SELECT, INSERT, UPDATE, DELETE ON TABLES TO shop_app;
ALTER DEFAULT PRIVILEGES FOR ROLE shop_owner IN SCHEMA public
    GRANT USAGE, SELECT ON SEQUENCES TO shop_app;
\end{sourcecode}

Die Auslassung ist \code{\textbackslash{}connect shop}: Ab dort läuft das Skript in der neuen
Datenbank. Wichtig sind die \code{DEFAULT PRIVILEGES}: Jede Tabelle, die eine spätere Migration als
\code{shop\_\allowbreak{}owner} anlegt, ist für \code{shop\_\allowbreak{}app} sofort nutzbar.
Ohne sie scheitert die neue Anwendungsversion mit \code{permission denied}, obwohl Migration und
Policy stimmen. Dasselbe gilt für Sequenzen: Ein \code{INSERT} in eine Tabelle mit
\code{serial}-Spalte ruft \code{nextval} auf ihrer Sequenz auf, EFs HiLo-Verfahren tut das vor
dem \code{INSERT}; ohne \code{USAGE} auf der Sequenz erhält die App-Rolle \code{permission
denied for sequence}. Eine Identity-Spalte (\code{GENERATED ... AS IDENTITY}, der Standard von
Npgsql) braucht für ein \code{INSERT} kein eigenes Recht auf ihre interne Sequenz. Die
Operation \code{Grant\allowbreak{}Table\allowbreak{}Privileges} des Plugins setzt ebenfalls
Default Privileges für Tabellen und Sequenzen, aber nur für Objekte, die die ausführende Rolle
später anlegt; das Skript legt sie für \code{shop\_\allowbreak{}owner} fest, unabhängig davon,
wer es ausführt.

\subsection*{Die Pipeline}

\begin{flowdiagram}
  \flowstep{Build und Test}{Ein Commit, ein Build: Tests, Bundle und Anwendung entstehen
  gemeinsam}
  \flowstep{Migrations-Bundle}{\code{dotnet ef migrations bundle}, eigenständig lauffähig;
  dazu das idempotente Skript zur Durchsicht}
  \flowstep{Freigabe}{GitHub-Environment \code{production}: Freigabe und Secrets nur für die
  folgenden Jobs}
  \flowstep{Bundle als shop\_owner}{Wendet ausstehende Migrationen in Transaktionen an; der
  einzige Job mit der Verbindung des Eigentümers}
  \flowstep{Rolling Deployment}{Erst nach erfolgreicher Migration; die Anwendung verbindet
  sich als \code{shop\_\allowbreak{}app}}
  \flowstep{Contract (später)}{Eine Folgemigration räumt auf, wenn keine alte Instanz mehr
  läuft}
\end{flowdiagram}

Die Vorlage \code{docs/\allowbreak{}deployment/\allowbreak{}migrations-\allowbreak{}deploy.\allowbreak{}yml}
setzt das mit drei Jobs um. Sie ist im Repository nicht aktiv; eine Anwendung kopiert sie
nach \code{.github/\allowbreak{}workflows} und ersetzt den Deploy-Schritt durch den ihrer
Plattform:

% source: docs/deployment/migrations-deploy.yml
\begin{sourcecode}[docs/\allowbreak{}deployment/\allowbreak{}migrations-\allowbreak{}deploy.\allowbreak{}yml (\allowbreak{}gekürzt)]{text}
jobs:
  build:
    steps:
      ...
      - name: Migration bundle
        run: >
          dotnet ef migrations bundle --project "$APP_PROJECT" --context "$DB_CONTEXT"
          --configuration Release --self-contained -r linux-x64 -o artifacts/efbundle
      ...
      - run: dotnet publish "$APP_PROJECT" -c Release -o artifacts/app

  migrate:
    needs: build
    environment: production
    steps:
      ...
      - name: Apply pending migrations as shop_owner
        env:
          OWNER_CONNECTION: ${{ secrets.SHOP_OWNER_CONNECTION_STRING }}
        run: |
          chmod +x artifacts/efbundle            # artifacts lose the executable bit
          artifacts/efbundle --connection "$OWNER_CONNECTION"

  deploy:
    needs: migrate
    environment: production
    ...
\end{sourcecode}

Das Bundle enthält das Startprojekt mit dem Plugin; für dessen Optionen gilt, was
Abschnitt „Migration Bundles“ in Kapitel~\ref{ch:advanced-migrations} sagt.
\code{--connection} überschreibt die Verbindung der Design-Time-Factory, und ein zweiter Lauf
wendet nichts an. Der Verbindungsstring des Eigentümers ist nur dem Job \code{migrate} bekannt; den der
Anwendung erhält die Hosting-Plattform, nicht die Pipeline.

Der Consumer-Test (\code{eng/\allowbreak{}consumer-\allowbreak{}test/\allowbreak{}run.\allowbreak{}sh
--nativeaot}, Seite~\pageref{sec:consumer-test}) spielt genau diesen Ablauf gegen PostgreSQL in Docker durch: Er führt
\code{setup-\allowbreak{}roles.\allowbreak{}sql} unverändert aus, baut das Bundle aus der
erzeugten Migration, wendet es als \code{shop\_\allowbreak{}owner} zweimal an (der zweite Lauf
ändert nichts), prüft, dass \code{shop\_\allowbreak{}app} die Tabelle nicht ändern darf
(\code{must be owner}), und startet die Anwendung als \code{shop\_\allowbreak{}app}, die nur
die Zeilen ihres Mandanten sieht.

\subsection*{Zero-Downtime unter RLS: Erweitern und Zusammenziehen}

Bei einem Rolling oder Blue-Green Deployment läuft die Migration vor der neuen Version, und
eine Zeit lang arbeiten alte und neue Instanzen gegen dasselbe Schema. Jede Migration muss
deshalb mit der Vorgängerversion der Anwendung funktionieren. Das Muster Expand and Contract
teilt eine Änderung dafür in drei Phasen, wie schon die Umstellung in
Abschnitt~„Von der Tabelle ohne Mandant zu RLS“:

\begin{datatable}{H{0.70}H{1.15}H{1.15}}
  \textbf{Phase} & \textbf{Schema} & \textbf{Unter RLS beachten} \\
  \tablehead
  1 Erweitern & Neue Spalten nullable oder mit Default, neue Tabellen, neue Indizes; nichts
  wird entfernt oder umbenannt. & Neue Tabellen mit \code{ITenant\allowbreak{}Owned} bekommen
  ihre Policy sofort; die alte Version benutzt sie nicht. Muss eine Policy sich ändern, kommt
  eine zweite hinzu, die die neue Form erlaubt. \\
  \tablerule
  2 Ausrollen & Unverändert. Die neue Version schreibt in die neuen Spalten, alte Instanzen
  laufen aus. & Beide Versionen setzen denselben Mandanten; die Policies müssen die Zeilen
  beider Versionen zulassen. \\
  \tablerule
  3 Zusammenziehen & Folgemigration: alte Spalten löschen, \code{NOT NULL} setzen, temporäre
  Policies entfernen. & Erst wenn keine alte Instanz mehr läuft; vorher ein Wächter wie in
  Abschnitt~„Von der Tabelle ohne Mandant zu RLS“. \\
\end{datatable}

Für Policies gilt eine Eigenheit von PostgreSQL: Mehrere \code{PERMISSIVE}-Policies einer
Tabelle werden mit \code{OR} verknüpft. Eine zusätzliche Policy erweitert also, was sichtbar
ist, und eignet sich als Übergang in Phase~1; ihr Entfernen in Phase~3 ist das
Zusammenziehen. Jede dieser Policies muss für sich mandantentrennend sein, sonst öffnet der
Übergang die Tabelle. Weil DDL in PostgreSQL transaktional ist, gibt es beim Ersetzen einer
Policy innerhalb einer Migration keinen Moment ohne Policy.

Ein Rollback der Anwendung ist bis Phase~3 ohne Datenbankänderung möglich, weil das Schema die
alte Version noch trägt. Danach braucht er das Down der Folgemigration, und gelöschte Spalten
kommen damit nicht zurück. Ein Down gehört deshalb nicht in die Pipeline; ein Fehler nach
Phase~3 wird mit einer neuen Migration vorwärts behoben.

% ------------------------------------------------------------
\chapter{Performance von Plugins}\label{ch:plugin-performance}
\versionrisk{stable}

Plugin-Code läuft in drei Häufigkeiten. Wer weiß, welcher Code wie oft läuft, weiß, wo sich
Optimierung lohnt:

\begin{datatable}{H{0.75}H{1.25}H{1.00}}
  \textbf{Häufigkeit} & \textbf{Code} & \textbf{Folgerung} \\
  \tablehead
  Einmal pro internem Provider oder Modell & Conventions, Type-Mapping-Plugins, Singleton-Options & Lesbarkeit
  vor Mikro-Optimierung \\
  \tablerule
  Einmal pro Abfrageform & Translators, Query-Interceptors & Moderat; Ergebnis wird gecacht
  (Kapitel~\ref{ch:the-query-pipeline}) \\
  \tablerule
  Bei jeder Ausführung & Interceptors für Commands, Verbindungen, SaveChanges,
  Materialisierung; Evaluatable-Filter, wenn der Funcletizer die Parameterwerte ausliest & Heißer Pfad: keine Reflection, keine unnötigen Allokationen \\
\end{datatable}

Daraus ergeben sich einige Regeln, die das Beispiel befolgt: \code{Method\allowbreak{}Info}-
und \code{Property\allowbreak{}Info}-Lookups liegen in statischen Feldern. Der
SaveChanges-Interceptor legt seine Liste erst an, wenn er tatsächlich eine mandantenfähige
Änderung findet. Der Connection-Interceptor kostet einen Roundtrip pro Öffnen; bei vielen
kleinen Abfragen in einer Anfrage öffnet die Anwendung die Verbindung deshalb einmal mit
\code{Open\allowbreak{}Connection\allowbreak{}Async(\allowbreak{})}, und es gilt „Ein Context,
ein Mandant“ (Kapitel~\ref{ch:interceptors}). Solche Annahmen sollte man messen, nicht raten. Das Repository enthält ein
BenchmarkDotNet-Projekt; der Benchmark für die SaveChanges-Logik läuft ohne Datenbank:

% source: benchmarks/Acme.EntityFrameworkCore.Tenancy.Benchmarks/SaveChangesInterceptorBenchmark.cs
\begin{sourcecode}[benchmarks/\allowbreak{}.\allowbreak{}.\allowbreak{}.\allowbreak{}/\allowbreak{}Save\allowbreak{}Changes\allowbreak{}Interceptor\allowbreak{}Benchmark.\allowbreak{}cs (\allowbreak{}gekürzt)]{csharp}
[MemoryDiagnoser]
public class SaveChangesInterceptorBenchmark
{
    [Params(1, 100)]
    public int Entities { get; set; }

    // GlobalSetup: Context mit UseAcmeTenancy() (ohne Verbindung) und dem Interceptor

    [Benchmark(Baseline = true)]
    public int TrackOnly()                  // was EF vor dem Speichern ohnehin tut
    {
        AddOrders();
        var count = _db.ChangeTracker.Entries().Count();
        _db.ChangeTracker.Clear();
        return count;
    }

    [Benchmark]
    public int TrackAndApplyTenant()        // dasselbe plus die SaveChanges-Logik des Plugins
    {
        AddOrders();
        _interceptor.ApplyTenant(_db);
        var count = _db.ChangeTracker.Entries().Count();
        _db.ChangeTracker.Clear();
        return count;
    }
}
\end{sourcecode}

\begin{sourcecode}[Benchmark ausführen]{text}
dotnet run -c Release --project benchmarks/Acme.EntityFrameworkCore.Tenancy.Benchmarks \
    -- --filter '*SaveChanges*' --job short
\end{sourcecode}

\begin{datatable}{H{1.67}H{0.83}H{0.83}H{0.83}H{0.83}}
  \textbf{Methode} & \textbf{Entitäten} & \textbf{Mittelwert} & \textbf{Verhältnis} & \textbf{Allokiert} \\
  \tablehead
  \code{Track\allowbreak{}Only} & 1 & 1,5 µs & 1,00 & 1,8 KB \\
  \tablerule
  \code{Track\allowbreak{}And\allowbreak{}Apply\allowbreak{}Tenant} & 1 & 2,3 µs & 1,54 &
  2,7 KB \\
  \tablerule
  \code{Track\allowbreak{}Only} & 100 & 123 µs & 1,00 & 134 KB \\
  \tablerule
  \code{Track\allowbreak{}And\allowbreak{}Apply\allowbreak{}Tenant} & 100 & 170 µs & 1,38 &
  202 KB \\
\end{datatable}

Beispielergebnisse (Short Job, .NET 10, Rechner des Autors; die absoluten Werte schwanken,
aussagekräftiger sind die Verhältnisse). Das Plugin fügt pro gespeicherter Entität etwa eine
halbe Mikrosekunde hinzu, rund 40 bis 50 Prozent mehr als das Change Tracking allein. Das
klingt nach viel, doch ein einziger Roundtrip zur Datenbank kostet Hunderte Mikrosekunden;
der zusätzliche Roundtrip des Connection-Interceptors pro Öffnen
(Kapitel~\ref{ch:interceptors}) wiegt daher weit schwerer als die SaveChanges-Logik. Genau
diese Art Erkenntnis liefert ein Benchmark: Er zeigt, wo sich Optimieren lohnt und wo nicht.

\section{Lesen und Abfragekompilierung}

Zwei weitere Benchmarks decken die anderen Häufigkeiten ab.
\code{Materialization\allowbreak{}Interceptor\allowbreak{}Benchmark} misst, was
\code{Verify\allowbreak{}Materialized\allowbreak{}Tenants(\allowbreak{})} beim Lesen kostet. Es
braucht echte Zeilen aus einem \code{Db\allowbreak{}Data\allowbreak{}Reader} und startet dafür
PostgreSQL per Testcontainers (Docker). Die Verbindung bleibt offen, sodass der Roundtrip des
Connection-Interceptors nicht mitgemessen wird; \code{As\allowbreak{}No\allowbreak{}Tracking}
hält das Change Tracking heraus. Eine dritte Variante mit einem leeren
\code{IMaterialization\allowbreak{}Interceptor} trennt den Aufwand von EF vom Aufwand des
Plugins:

% source: benchmarks/Acme.EntityFrameworkCore.Tenancy.Benchmarks/MaterializationInterceptorBenchmark.cs
\begin{sourcecode}[benchmarks/\allowbreak{}.\allowbreak{}.\allowbreak{}.\allowbreak{}/\allowbreak{}Materialization\allowbreak{}Interceptor\allowbreak{}Benchmark.\allowbreak{}cs (gekürzt)]{csharp}
[MemoryDiagnoser]
public class MaterializationInterceptorBenchmark
{
    [Params(10, 100, 1000)]
    public int Rows { get; set; }

    // GlobalSetup: Container, 1000 Zeilen eines Mandanten, Verbindungen offen

    [Benchmark(Baseline = true)]
    public Task<int> Read() => CountAsync(_plain);

    [Benchmark]
    public Task<int> ReadIntercepted() => CountAsync(_intercepted);

    [Benchmark]
    public Task<int> ReadAndVerify() => CountAsync(_verifying);

    private async Task<int> CountAsync(BenchmarkContext db)
        => (await db.Orders.AsNoTracking().OrderBy(o => o.Id).Take(Rows).ToListAsync()).Count;

    // Create: UseAcmeTenancy(sp, t => t.VerifyMaterializedTenants(verify))
}
\end{sourcecode}

\begin{datatable}{H{1.67}H{0.83}H{0.83}H{0.83}H{0.83}}
  \textbf{Methode} & \textbf{Zeilen} & \textbf{Mittelwert} & \textbf{Verhältnis} & \textbf{Allokiert} \\
  \tablehead
  \code{Read} & 10 & 200 µs & 1,00 & 12 KB \\
  \tablerule
  \code{Read\allowbreak{}Intercepted} & 10 & 210 µs & 1,08 & 21 KB \\
  \tablerule
  \code{Read\allowbreak{}And\allowbreak{}Verify} & 10 & 220 µs & 1,12 & 22 KB \\
  \tablerule
  \code{Read} & 100 & 278 µs & 1,00 & 36 KB \\
  \tablerule
  \code{Read\allowbreak{}Intercepted} & 100 & 764 µs & 2,77 & 121 KB \\
  \tablerule
  \code{Read\allowbreak{}And\allowbreak{}Verify} & 100 & 741 µs & 2,68 & 136 KB \\
  \tablerule
  \code{Read} & 1000 & 846 µs & 1,00 & 269 KB \\
  \tablerule
  \code{Read\allowbreak{}Intercepted} & 1000 & 3858 µs & 4,70 & 1120 KB \\
  \tablerule
  \code{Read\allowbreak{}And\allowbreak{}Verify} & 1000 & 4288 µs & 5,22 & 1268 KB \\
\end{datatable}

Das Ergebnis überrascht: Den größten Teil kostet nicht die Prüfung, sondern dass überhaupt ein
Materialisierungs-Interceptor registriert ist. EF erzeugt dann für jede Zeile die Daten, die
\code{Materialization\allowbreak{}Interception\allowbreak{}Data} anbietet, und verlässt den
schnellen, kompilierten Materialisierungspfad: bei 1000 Zeilen etwa 0,85 KB und 3 µs pro
Zeile. Die Prüfung des Plugins (Provider aus dem Cache, Wert lesen, vergleichen) fügt rund
0,15 KB und 0,4 µs pro Zeile hinzu. Bei zehn Zeilen geht beides im Roundtrip unter, bei
tausend Zeilen wird das Lesen auf dem Client etwa fünfmal so teuer.

\code{Query\allowbreak{}Compilation\allowbreak{}Benchmark} braucht keine Datenbank:
\code{To\allowbreak{}Query\allowbreak{}String} durchläuft die Pipeline bis zum SQL ohne
Verbindung. \code{Compile} leert vor jedem Aufruf den Query-Cache von EF (den
\code{IMemory\allowbreak{}Cache} des internen Service Providers), \code{Cache\allowbreak{}Hit}
führt dieselbe Abfrageform mit neuem Parameterwert aus. Verglichen werden EF Core allein, das
Plugin und das Plugin mit Query-Filter:

\begin{datatable}{H{1.40}H{0.90}H{0.90}H{0.90}H{0.90}}
  \textbf{Context} & \textbf{Kompilieren} & \textbf{Allokiert} & \textbf{Cache-Treffer} &
  \textbf{Allokiert} \\
  \tablehead
  EF Core & 2,64 ms & 146 KB & 18 µs & 9,9 KB \\
  \tablerule
  Plugin & 2,68 ms & 147 KB & 15 µs & 9,9 KB \\
  \tablerule
  Plugin mit Query-Filter & 2,78 ms & 160 KB & 15 µs & 10,2 KB \\
\end{datatable}

Translator-Plugins und Evaluatable-Filter kosten beim Kompilieren nichts Messbares; der Query-Filter verlängert den Ausdrucksbaum und damit die
Übersetzung um wenige Prozent. Bei einem Cache-Treffer verschwindet jeder Unterschied im
Rauschen. Alle Werte stammen aus einem Short Job (.NET 10, Rechner des Autors); dessen
Fehlerbalken sind breit, die Größenordnungen aber eindeutig:

\begin{datatable}{H{0.80}H{0.90}H{1.30}}
  \textbf{Phase} & \textbf{Häufigkeit} & \textbf{Mehraufwand des Plugins} \\
  \tablehead
  Kompilieren & Einmal pro Abfrageform & Kaum messbar, mit Query-Filter wenige Prozent \\
  \tablerule
  Cache-Treffer & Jede Abfrage & Nicht messbar \\
  \tablerule
  SaveChanges & Jedes Speichern & Etwa 0,5 µs pro Entität (40 bis 50 Prozent des Change
  Trackings) \\
  \tablerule
  Materialisierung & Jede gelesene Zeile & Etwa 3,5 µs pro Zeile mit
  \code{Verify\allowbreak{}Materialized\allowbreak{}Tenants(\allowbreak{})}, ohne: keiner \\
\end{datatable}

Ob alle Benchmarks laufen, zeigt ein Probelauf ohne aussagekräftige Messwerte:

\begin{sourcecode}[Alle Benchmarks, Probelauf]{text}
dotnet run -c Release --project benchmarks/Acme.EntityFrameworkCore.Tenancy.Benchmarks \
    -- --filter '*' --job dry
\end{sourcecode}

\subsection*{Wann die Prüfung beim Lesen lohnt}

\code{Verify\allowbreak{}Materialized\allowbreak{}Tenants(\allowbreak{})} ist Defense in Depth:
Es schützt nicht vor einem Angriff, den RLS nicht abwehrt, sondern bemerkt, wenn RLS gar nicht
greift, etwa weil die Anwendung mit dem Eigentümer, einem Superuser oder einer Rolle mit
\code{BYPASSRLS} verbindet oder eine selbst geöffnete Verbindung ohne
\code{Apply\allowbreak{}Tenant\allowbreak{}Session} arbeitet. Daraus ergibt sich die
Empfehlung:

\begin{itemize}
  \item \textbf{Einschalten}, wenn eine Fehlkonfiguration teurer ist als die Rechenzeit:
    Daten mit hohen Schutzanforderungen, viele Umgebungen mit eigener Konfiguration, kleine
    Ergebnismengen, wie sie typische API-Abfragen liefern. Dort geht der Aufwand im Roundtrip
    unter.
  \item \textbf{Ausschalten} für leselastige Pfade mit großen Ergebnismengen, etwa Exporte und
    Berichte. Dort trägt die Datenbank die Isolation allein, abgesichert durch die
    Rollentrennung (Seite~\pageref{sec:deployment}) und Tests, die mit der eingeschränkten
    Rolle der Anwendung lesen.
  \item \textbf{Dazwischen}: in Tests und Staging immer einschalten, in Produktion nur für den
    Context der sensiblen Daten. Die Option gilt pro Context-Konfiguration, nicht global.
\end{itemize}

\begin{tip}
  Für den Start großer Modelle ist das Compiled Model
  (Kapitel~\ref{ch:compiled-models-trimming-and-nativeaot}) der wirksamste Hebel.
  Conventions eines Plugins laufen dann gar nicht mehr zur Laufzeit; deshalb lohnt es sich,
  die Compiled-Model-Fähigkeit zu testen statt nur zu behaupten.
\end{tip}

\section{Second-Level-Caches}

Ein Second-Level-Cache (etwa \code{EFCore\allowbreak{}Second\allowbreak{}Level\allowbreak{}Cache\allowbreak{}Interceptor}
oder Ergebnisse in Redis oder \code{IMemory\allowbreak{}Cache}) ist der naheliegende nächste
Schritt, wenn Abfragen zu langsam sind. Mit Mandantenfähigkeit ist er auch der kürzeste Weg zu
einem Datenleck. Der Cache erkennt ein Ergebnis am SQL-Text oder am normalisierten
LINQ-Ausdruck, und bei Row Level Security steht der Mandant in keinem von beiden: Die Policy
liest ihn aus der Session.

\begin{datatable}{H{0.85}H{1.15}}
  \textbf{Schritt} & \textbf{Was passiert} \\
  \tablehead
  Mandant A fragt ab & Die Datenbank filtert nach dem Session-Wert, der Cache speichert die
  Zeilen von A unter dem SQL-Text \\
  \tablerule
  Mandant B stellt dieselbe Abfrage & Gleicher SQL-Text, gleicher Schlüssel: Der Cache liefert
  die Zeilen von A. Die Datenbank wird nie gefragt, die Policies laufen nie \\
\end{datatable}

Die übrigen Schutzmechanismen des Plugins helfen hier auch nicht. Der SaveChanges-Interceptor
sieht nur das Speichern. Die Prüfung beim Laden
(Kapitel~\ref{ch:materialization-from-rows-to-objects}) läuft nur, wenn EF Entitäten aus Zeilen
baut; ein Cache, der Entitäten oder DTOs speichert, gibt sie ohne diesen Schritt heraus. Ein
Filter mit dem Mandanten als Abfrageparameter (Cosmos DB, Testmodus) bringt den Mandanten nur
dann in den Schlüssel, wenn der Cache Parameterwerte einbezieht, und das tut nicht jeder.

\begin{important}[Der Mandant gehört in jeden Cache-Schlüssel]
  Ein Cache vor einer mandantenbezogenen Abfrage ist nur sicher, wenn sein Schlüssel den
  aktuellen Mandanten enthält, gesetzt von der Anwendung und nicht aus dem SQL abgeleitet. Ohne
  Mandant wird nichts gecacht. Die meisten Cache-Bibliotheken haben einen Haken für den
  Schlüssel oder ein Schlüsselpräfix; wo er fehlt, gehört der Cache nicht vor Mandantendaten.
  Ein Test mit zwei Mandanten und derselben Abfrage zeigt, ob der Haken wirkt.
\end{important}

\begin{sourcecode}[Beispiel]{csharp}
// Überall, wo Abfrageergebnisse gecacht werden: Der Mandant gehört in jeden Schlüssel
public sealed class TenantCacheKeys
{
    private readonly ITenantProvider _tenants;

    public TenantCacheKeys(ITenantProvider tenants) => _tenants = tenants;

    public string For(string queryKey)
        => _tenants.CurrentTenantId is { } tenant
            ? $"tenant:{tenant}:{queryKey}"
            : throw new InvalidOperationException("No tenant: do not cache the result.");
}
\end{sourcecode}

% ------------------------------------------------------------
\chapter{Einsatz in ASP.NET Core}\label{ch:using-the-plugin-in-asp}
\versionrisk{stable}

Das Beispiel \code{samples/\allowbreak{}Shop\allowbreak{}Api} zeigt den Weg vom Plugin in
eine Web-API. Der Mandant kommt aus dem Claim \code{tenant\_\allowbreak{}id} des
angemeldeten Benutzers:

% source: samples/ShopApi/HttpTenantProvider.cs
\begin{sourcecode}[samples/\allowbreak{}Shop\allowbreak{}Api/\allowbreak{}Http\allowbreak{}Tenant\allowbreak{}Provider.\allowbreak{}cs]{csharp}
public sealed class HttpTenantProvider : ITenantProvider
{
    public const string ClaimType = "tenant_id";

    private readonly IHttpContextAccessor _httpContextAccessor;

    public HttpTenantProvider(IHttpContextAccessor httpContextAccessor)
    {
        ArgumentNullException.ThrowIfNull(httpContextAccessor);

        _httpContextAccessor = httpContextAccessor;
    }

    // Bei jedem Zugriff neu ermittelt: funktioniert auch als Singleton (Pooling)
    public TenantId? CurrentTenantId
    {
        get
        {
            var value = _httpContextAccessor.HttpContext?.User.FindFirst(ClaimType)?.Value;
            return Guid.TryParse(value, out var id) ? new TenantId(id) : null;
        }
    }
}
\end{sourcecode}

% source: samples/ShopApi/Program.cs
\begin{sourcecode}[samples/\allowbreak{}Shop\allowbreak{}Api/\allowbreak{}Program.\allowbreak{}cs (\allowbreak{}Auszug)]{csharp}
builder.Services.AddHttpContextAccessor();
builder.Services.AddAcmeTenancy<HttpTenantProvider>();
builder.Services.AddDbContext<ShopApiContext>((sp, options) => options
    .UseNpgsql(builder.Configuration.GetConnectionString("Shop"))
    .UseAcmeTenancy(sp));

// Kein Mandantenfilter im Code: RLS in der Datenbank sorgt für die Trennung
app.MapGet("/orders", async (ShopApiContext db) =>
    await db.Orders.OrderBy(o => o.Id).ToListAsync());
\end{sourcecode}

In der Entwicklungsumgebung legt das Beispiel beim Start seine Datenbank, die Tabellen mit RLS
und die App-Rolle selbst an, und der Header \code{X-\allowbreak{}Demo-\allowbreak{}Tenant}
simuliert den Claim. Mit der Datenbank aus \code{docker-compose.\allowbreak{}yml} läuft das
ganze Beispiel so:

% run: shopapi
\begin{sourcecode}[Shell]{text}
docker compose up -d
dotnet run --project samples/ShopApi

curl -X POST -H "X-Demo-Tenant: 0199f0d2-0000-7000-8000-00000000000a" \
    "http://localhost:5000/orders?total=10"
curl -X POST -H "X-Demo-Tenant: 0199f0d2-0000-7000-8000-00000000000b" \
    "http://localhost:5000/orders?total=99"
curl -H "X-Demo-Tenant: 0199f0d2-0000-7000-8000-00000000000a" \
    "http://localhost:5000/orders"        # nur die Bestellung von Mandant ...0a
curl "http://localhost:5000/orders"       # ohne Mandant: []
\end{sourcecode}

\begin{important}[Woher der Mandant kommt]
  Der Mandant muss aus einer vertrauenswürdigen Quelle kommen: einem Claim aus einem
  signierten Token. Ein Header oder Query-Parameter, den der Client frei setzt, hebelt die
  gesamte Isolation aus. Das Beispiel enthält eine Header-Simulation nur für die
  Entwicklungsumgebung und kennzeichnet sie ausdrücklich.
\end{important}

Für \code{Add\allowbreak{}Db\allowbreak{}Context\allowbreak{}Pool} registriert man den
Provider als Singleton. Das ist möglich, weil
\code{Http\allowbreak{}Tenant\allowbreak{}Provider} den Mandanten bei jedem Zugriff neu
ermittelt; die Interceptors fragen ihn bei jeder Verbindung und jedem Speichern:

\begin{sourcecode}[Variante mit Pooling]{csharp}
builder.Services.AddAcmeTenancy<HttpTenantProvider>(ServiceLifetime.Singleton);
builder.Services.AddDbContextPool<ShopApiContext>((sp, options) => options
    .UseNpgsql(connectionString)
    .UseAcmeTenancy(sp));   // sp: Root-Provider; Interceptors entstehen einmal
\end{sourcecode}

\section{DbContext-Factories}

\code{Add\allowbreak{}Db\allowbreak{}Context\allowbreak{}Factory} ist die übliche Wahl in
Blazor Server und für parallele Arbeit. Anders als
\code{Add\allowbreak{}Db\allowbreak{}Context} registriert es Factory und Options
standardmäßig als \textbf{Singletons}. Der Options-Callback bekommt dann den Root-Provider,
nicht den Scope der aktuellen Anfrage. Für einen Scoped-Tenant-Provider, der Zustand pro
Scope hält, hat das Folgen:

\begin{datatable}{H{0.87}H{1.13}}
  \textbf{Registrierung} & \textbf{Verhalten} \\
  \tablehead
  Singleton-Factory, Scope-Validierung an (ASP.NET Core in Development) & Exception: „Cannot
  resolve scoped service 'ITenantProvider' from root provider.“ \\
  \tablerule
  Singleton-Factory, Scope-Validierung aus (Standard in Produktion) & Kein Fehler. Die
  Interceptors bekommen eine Instanz aus dem Root-Container, die keine Anfrage je befüllt:
  kein Mandant, Speichern wirft, Abfragen sehen nichts \\
  \tablerule
  Factory als \code{Scoped} registriert & Jeder Scope bekommt eine eigene Factory mit
  eigenem Tenant-Provider: korrekt \\
\end{datatable}

\begin{sourcecode}[Factory mit dem Mandanten der Anfrage]{csharp}
builder.Services.AddDbContextFactory<ShopApiContext>(
    (sp, options) => options.UseNpgsql(connectionString).UseAcmeTenancy(sp),
    ServiceLifetime.Scoped);         // Factory und Options pro Scope
\end{sourcecode}

Ein Provider, der den Mandanten bei jedem Zugriff neu ermittelt, etwa
\code{Http\allowbreak{}Tenant\allowbreak{}Provider}, funktioniert in jeder Variante; das
Problem betrifft Provider mit Zustand pro Scope.
\code{Add\allowbreak{}Pooled\allowbreak{}Db\allowbreak{}Context\allowbreak{}Factory} verhält
sich immer wie Pooling: Die Options existieren einmal, der Tenant-Provider muss also ein
solches Singleton sein. \code{Db\allowbreak{}Context\allowbreak{}Factory\allowbreak{}Tests}
zeigt alle drei Varianten.

\section{Außerhalb von HTTP-Anfragen}

Hintergrunddienste, Message-Handler und geplante Jobs haben keine Anfrage und keinen Claim.
Der Mandant kommt dann aus dem Job selbst, etwa aus der Nachricht oder aus einer Liste zu
verarbeitender Mandanten. Der \code{Ambient\allowbreak{}Tenant\allowbreak{}Provider} des
Beispiels speichert ihn in einem \code{Async\allowbreak{}Local}, das mit dem asynchronen
Kontrollfluss mitwandert:

% source: samples/ShopSample/AmbientTenantProvider.cs
\begin{sourcecode}[samples/\allowbreak{}Shop\allowbreak{}Sample/\allowbreak{}Ambient\allowbreak{}Tenant\allowbreak{}Provider.\allowbreak{}cs (\allowbreak{}gekürzt)]{csharp}
public sealed class AmbientTenantProvider : ITenantProvider
{
    private static readonly AsyncLocal<TenantId?> Current = new();

    public TenantId? CurrentTenantId => Current.Value;

    public static IDisposable Use(TenantId tenant)
    {
        var previous = Current.Value;
        Current.Value = tenant;
        return new Restore(previous);     // setzt beim Dispose den vorigen Mandanten zurück
    }

    // private sealed class Restore : IDisposable { ... Current.Value = _previous; }
}
\end{sourcecode}

\begin{sourcecode}[Job-Muster (\allowbreak{}tests/\allowbreak{}.\allowbreak{}.\allowbreak{}.\allowbreak{}/\allowbreak{}Background\allowbreak{}Tenant\allowbreak{}Tests.\allowbreak{}cs,\allowbreak{} vereinfacht)]{csharp}
// Erst der Mandant, dann Scope und Context: die Interceptors entstehen im Scope
using (AmbientTenantProvider.Use(tenant))
{
    await using var scope = scopeFactory.CreateAsyncScope();
    var db = scope.ServiceProvider.GetRequiredService<ShopContext>();
    // ... Arbeit für diesen Mandanten ...
}
\end{sourcecode}

Weil der Provider das \code{Async\allowbreak{}Local} bei jedem Zugriff liest, kann er als
Singleton registriert werden. Parallele Jobs behalten ihren eigenen Mandanten, und nach dem
\code{using}-Block bleibt kein Mandant zurück;
\code{Background\allowbreak{}Tenant\allowbreak{}Tests} prüft beides mit 20 parallelen Jobs.
Eine Anwendung mit Web-Anfragen \emph{und} Jobs kombiniert beide Quellen in einem Provider:
zuerst der Claim, sonst der Ambient-Mandant.

\textbf{Nachrichten.} Bei einem Message Broker (RabbitMQ, Kafka, Azure Service Bus) oder einem
Framework darauf (MassTransit, Wolverine) schreibt der Sender den aktuellen Mandanten in einen
Header der Nachricht, und die Empfängerseite setzt ihn, bevor der Consumer läuft. Der Ort dafür
ist der Consume-Filter oder die Middleware des Frameworks, damit kein einzelner Consumer es
vergessen kann:

\begin{sourcecode}[Consume-Filter (\allowbreak{}vereinfacht,\allowbreak{} unabhängig vom Framework)]{csharp}
// Wo der Broker eine Nachricht übergibt: Mandant aus dem Header, dann Scope und Context
public async Task HandleAsync(BrokerMessage message, CancellationToken cancellationToken)
{
    if (!message.Headers.TryGetValue("X-Tenant-Id", out var value)
        || !Guid.TryParse(value, out var id))
        throw new InvalidOperationException("Message without a tenant: rejected.");

    using (AmbientTenantProvider.Use(new TenantId(id)))
    {
        await using var scope = scopeFactory.CreateAsyncScope();
        var handler = scope.ServiceProvider.GetRequiredService<IOrderHandler>();
        await handler.HandleAsync(message, cancellationToken);   // nutzt ShopContext
    }
}
\end{sourcecode}

Zwei Regeln übertragen sich von den Jobs. Der Header wird gelesen, bevor der Scope entsteht,
denn Context und Interceptors leben im Scope. Und eine Nachricht ohne gültigen Mandanten wird
abgewiesen und landet in der Dead-Letter-Queue; sie wird nie für alle Mandanten verarbeitet.
Ein Scoped-Tenant-Provider mit Setter funktioniert genauso, solange der Filter den Mandanten im
Scope des Consumers setzt, bevor der Context aufgelöst wird.

\begin{important}[Nie auf „alle Mandanten“ zurückfallen]
  Ohne Mandant muss ein Job scheitern oder nichts sehen, niemals alles. Das Plugin verhält
  sich so: Der Session-Wert bleibt leer, die Policies treffen keine Zeile, und Speichern
  wirft. Jobs, die wirklich alle Mandanten brauchen, nutzen die Reporting-Rolle aus
  Kapitel~\ref{ch:row-level-security-in-production}.
\end{important}

% ------------------------------------------------------------
\chapter{CI und API-Stabilität}\label{ch:ci-and-api-stability}
\versionrisk{stable}

\section{Der Workflow}

Der GitHub-Actions-Workflow im Repository setzt die Empfehlungen aus
Kapitel~\ref{ch:test-strategy} und \ref{ch:versioning-and-compatibility} um. Der Job
\code{build} baut und testet ohne und mit Datenbank, \code{handbook} baut das Handbuch,
\code{consumer} prüft die Pakete als Verbraucher. Der Job \code{cosmos} (Emulator, noch nicht
blockierend) fehlt im Auszug:

\begin{sourcecode}[.\allowbreak{}github/\allowbreak{}workflows/\allowbreak{}ci.\allowbreak{}yml (\allowbreak{}gekürzt)]{text}
jobs:
  build:
    runs-on: ubuntu-latest
    steps:
      - uses: actions/checkout@v4
      - uses: actions/setup-dotnet@v4
        with: { global-json-file: global.json }
      - run: dotnet build -c Release
      - run: dotnet test tests/Acme.EntityFrameworkCore.Tenancy.Tests -c Release --no-build
      - run: dotnet test tests/Acme.EntityFrameworkCore.Tenancy.IntegrationTests
          -c Release --no-build
      - run: dotnet test tests/Acme.EntityFrameworkCore.Tenancy.SqlServer.Tests
          -c Release --no-build

  handbook:                            # beide PDFs mit LuaLaTeX, im Ordner docs/handbook
    runs-on: ubuntu-latest
    container: texlive/texlive@sha256:540d…   # TeX Live 2026, per Digest fixiert
    steps:
      # ... checkout ...
      - run: python3 check_code.py       # Zeilenlänge, Auszüge gegen die Quellen
      - run: python3 check_structure.py  # beide Fassungen gleich, Verweise, Tabellen
      - run: latexmk EF-Core-Plugin-Handbuch.tex
      - run: latexmk EF-Core-Plugin-Handbook.tex

  consumer:                            # die Pakete als Verbraucher prüfen
    steps:
      - run: eng/consumer-test/run.sh --nativeaot

  handbook-examples:                   # die mit "% run:" markierten Listings ausführen
    steps:
      - run: python3 docs/handbook/check_examples.py
\end{sourcecode}

Codeauszüge veralten leise, wenn sich der Code ändert. Deshalb bindet eine Kommentarzeile
\code{\% source: <Pfad>} direkt vor dem Listing einen Auszug an seine Quelldatei, und
\code{check\_\allowbreak{}code.\allowbreak{}py} vergleicht beide: Kommentare und
Auslassungen (\code{...}) zählen nicht, sodass die deutsche Fassung Kommentare übersetzen und
Auszüge kürzen darf; jede übrige Zeile muss in derselben Reihenfolge in der Quelle stehen,
Zeilenumbrüche dürfen abweichen. Mit dem Zusatz \code{(strings)} bleiben auch übersetzte
Meldungen außen vor. Ein Auszug, dessen Titel eine Datei des Repositorys nennt und der nicht
gebunden ist, muss „vereinfacht“ heißen. \code{check\_\allowbreak{}structure.\allowbreak{}py}
prüft, dass beide Fassungen pro Kapitel dieselben Abschnitte, Listings, Bindungen, Tabellen und
Hinweisboxen haben, dass jedes Verweisziel existiert und die Spaltenbreiten jeder Tabelle
aufgehen.

Beide Skripte lesen das Buch; keines führt etwas aus. Ein Befehl, der erfolgreich endet, ohne
seine Aufgabe zu erfüllen, etwa ein Testfilter, der keinen Test trifft, käme durch. Einige
Listings werden deshalb \emph{ausgeführt}: Eine Kommentarzeile \code{\% run: <Art>} markiert
sie, und \code{check\_\allowbreak{}examples.\allowbreak{}py} führt sie im CI-Job
\code{handbook-\allowbreak{}examples} aus, in temporären Verzeichnissen und Containern, die es
wieder entfernt. Derzeit sind das zwei Listings: der Testbefehl im Übungsanhang, für den das
Skript die Zähler aus einer TRX-Ergebnisdatei liest und mindestens einen ausgeführten und keinen
fehlgeschlagenen Test verlangt, und die Sitzung von \code{Shop\allowbreak{}Api} in
Kapitel~\ref{ch:using-the-plugin-in-asp}, deren curl-Anfragen gegen das Beispiel und PostgreSQL
laufen und die Isolation zeigen müssen. Alle übrigen Listings werden nicht ausgeführt:
C\#-Auszüge sind mit \code{\% source:} gebunden oder als vereinfacht gekennzeichnet, Workflow-
und Konfigurationsauszüge sind gekürzt, und Listings mit Auslassungen sind keine Programme. Die
Pakete als Ganzes, das \code{dotnet ef}-Tooling und NativeAOT deckt der Consumer-Test ab, nicht
dieses Skript.

Das Beispiel baut nur gegen stabile Versionen (.NET 10, EF Core 10). Ein zweiter Workflow,
\code{compatibility.\allowbreak{}yml}, läuft wöchentlich und testet gegen die neuesten
stabilen 10.0.x-Patches, denn die reguläre CI baut gegen den Anfang des Versionsbereichs. Die
Kanäle und eine Skizze für die nächste EF-Major-Version mit Previews und RCs beschreibt
Abschnitt „Patches und Previews in der CI“ in Kapitel~\ref{ch:versioning-and-compatibility}.

\section{Öffentliche API absichern}

Ein Plugin hat zwei öffentliche APIs: die für Anwender (Fluent API, Options, Typen) und die
für Erweiterer (virtuelle Methoden, Dependencies, öffentliche Dienste). Beide brechen leicht
unbemerkt. Der
\code{Microsoft.\allowbreak{}Code\allowbreak{}Analysis.\allowbreak{}Public\allowbreak{}Api\allowbreak{}Analyzers}
führt die öffentliche API in zwei Textdateien
(\code{Public\allowbreak{}API.\allowbreak{}Shipped.\allowbreak{}txt},
\code{Public\allowbreak{}API.\allowbreak{}Unshipped.\allowbreak{}txt}) und meldet jede
Abweichung. Jede API-Änderung wird so im Pull Request als Diff sichtbar.

\begin{note}[Einrichtung]
  Beim ersten Build meldet der Analyzer jede öffentliche API als RS0016, weil die Dateien
  leer sind. Der Code-Fix „Add all items … to the public API“ füllt sie, auf der
  Kommandozeile mit
  \code{dotnet format analyzers -{}-diagnostics RS0016 -{}-severity warn}. Beim Release
  wandert der Inhalt von Unshipped nach Shipped; ab dann ist das Entfernen einer
  ausgelieferten API ein sichtbarer Breaking Change.
\end{note}

\begin{important}[Warnungen reichen nicht]
  Solange RS0016 und RS0017 nur Warnungen sind, gehen sie in der Build-Ausgabe unter. Im
  Beispiel blieben die Dateien so bis Ausgabe~1.0 leer: 178 Warnungen bei einem vollständigen
  Build, keine bei inkrementellen, und niemand bemerkte es.
  \code{Directory.\allowbreak{}Build.\allowbreak{}props} macht beide deshalb zu Fehlern
  (\code{Warnings\allowbreak{}As\allowbreak{}Errors}); eine neue öffentliche API muss jetzt
  in \code{Public\allowbreak{}API.\allowbreak{}Unshipped.\allowbreak{}txt} eingetragen
  werden, sonst bricht der Build.
\end{important}
